% GENERATED FILE. Edit canonical lessons, then run npm run generate:books.
% canonical-source-hash: fnv1a64:9dce7ed4f697fde8
% canonical-lessons: FR-C382-enfreindre, FR-C382-infraction, FR-C382-transgresser, FR-C382-sanction, FR-C382-sanctionner

\chapter{Enfreindre les règles}
\label{ch:fr-enfreindre-les-regles}
\hlchaptermodality{382}

\begin{chapteropening}
\textbf{By the end of this chapter:} \emph{I can say that someone broke a rule: breaking the law, an offence, transgressing, a penalty and penalising someone.}

Complete the last lesson of chapter 382: I can say that someone broke a rule: breaking the law, an offence, transgressing, a penalty and penalising someone.
\end{chapteropening}

\section[enfreindre]{enfreindre — to break (a rule, the law)}

\label{lesson:FR-C382-enfreindre}



\begin{quote}
Before the new one: say the French for the tolerance, then the French for to demand.
\end{quote}

\subsection*{You'll want to know: enfreindre}
\textbf{enfreindre} — \textquotedblleft{}to break (a rule, the law)\textquotedblright{}.

\begin{quote}
Il a enfreint le code de la route.
\end{quote}

\begin{quote}
\emph{He broke the highway code.}
\end{quote}

Conjugated like \emph{peindre}: \emph{il enfreint}, \emph{ils enfreignent}.

\begin{cousinweb}
\emph{enfreindre} is usually traced to a Vulgar Latin \emph{infrangere} \textquotedblleft{}to break into\textquotedblright{}, reshaped from Latin \emph{infringere} \textquotedblleft{}to break\textquotedblright{} (Lewis \& Short), built on \emph{frangere} \textquotedblleft{}to break\textquotedblright{}.
\end{cousinweb}

\begin{grammarlens}[title={What you've built}]
One more word for telling, explaining and describing at B1.
\end{grammarlens}

\subsection*{Your turn}
\begin{itemize}
\raggedright
  \item \emph{Say it:} \emph{enfreindre}
  \item \emph{Say it:} \emph{enfreindre}, once more
  \item \emph{Say it:} a sentence of your own with \emph{enfreindre}
  \item \emph{Recall:} say the French for the tolerance, then the French for to demand, then say \emph{enfreindre} again
\end{itemize}

\begin{tcolorbox}[breakable,colback=teal!4,colframe=teal!35!black,title={Before you move on}]
What does enfreindre mean? (\textquotedblleft{}To break (a rule, the law)\textquotedblright{}.) Say it once more.
\end{tcolorbox}
\section[l'infraction]{l'infraction — the offence, the violation}

\label{lesson:FR-C382-infraction}



\begin{quote}
Before the new one: say the French for to demand, then the French for to break.
\end{quote}

\subsection*{You'll want to know: l'infraction}
\textbf{l'infraction} — \textquotedblleft{}the offence, the violation\textquotedblright{}.

\begin{quote}
Conduire sans permis est une infraction grave.
\end{quote}

\begin{quote}
\emph{Driving without a licence is a serious offence.}
\end{quote}

Feminine: \emph{une infraction au code de la route}.

\begin{cousinweb}
\emph{infraction} is a learned borrowing from Latin \emph{infractiō} \textquotedblleft{}a breaking\textquotedblright{} (EtymDB), built on \emph{frangere} \textquotedblleft{}to break\textquotedblright{} (Lewis \& Short).
\end{cousinweb}

\begin{grammarlens}[title={What you've built}]
One more word for telling, explaining and describing at B1.
\end{grammarlens}

\subsection*{Your turn}
\begin{itemize}
\raggedright
  \item \emph{Say it:} \emph{l'infraction}
  \item \emph{Say it:} \emph{l'infraction}, once more
  \item \emph{Say it:} a sentence of your own with \emph{l'infraction}
  \item \emph{Recall:} say the French for to demand, then the French for to break, then say \emph{l'infraction} again
\end{itemize}

\begin{tcolorbox}[breakable,colback=teal!4,colframe=teal!35!black,title={Before you move on}]
What does l'infraction mean? (\textquotedblleft{}The offence, the violation\textquotedblright{}.) Say it once more.
\end{tcolorbox}
\section[transgresser]{transgresser — to transgress, to break (a rule)}

\label{lesson:FR-C382-transgresser}



\begin{quote}
Before the new one: say the French for to break, then the French for the offence.
\end{quote}

\subsection*{You'll want to know: transgresser}
\textbf{transgresser} — \textquotedblleft{}to transgress, to break (a rule)\textquotedblright{}.

\begin{quote}
Les adolescents aiment transgresser les règles.
\end{quote}

\begin{quote}
\emph{Teenagers like to break the rules.}
\end{quote}

Formal, a little literary.

\begin{cousinweb}
\emph{transgresser} is built on Latin \emph{trānsgressus}, from \emph{trānsgredī} \textquotedblleft{}to step across\textquotedblright{} (Lewis \& Short), built on \emph{gradī} \textquotedblleft{}to step\textquotedblright{}.
\end{cousinweb}

\begin{grammarlens}[title={What you've built}]
One more word for telling, explaining and describing at B1.
\end{grammarlens}

\subsection*{Your turn}
\begin{itemize}
\raggedright
  \item \emph{Say it:} \emph{transgresser}
  \item \emph{Say it:} \emph{transgresser}, once more
  \item \emph{Say it:} a sentence of your own with \emph{transgresser}
  \item \emph{Recall:} say the French for to break, then the French for the offence, then say \emph{transgresser} again
\end{itemize}

\begin{tcolorbox}[breakable,colback=teal!4,colframe=teal!35!black,title={Before you move on}]
What does transgresser mean? (\textquotedblleft{}To transgress, to break (a rule)\textquotedblright{}.) Say it once more.
\end{tcolorbox}
\section[la sanction]{la sanction — the penalty, the sanction}

\label{lesson:FR-C382-sanction}



\begin{quote}
Before the new one: say the French for the offence, then the French for to transgress.
\end{quote}

\subsection*{You'll want to know: la sanction}
\textbf{la sanction} — \textquotedblleft{}the penalty, the sanction\textquotedblright{}.

\begin{quote}
Les sanctions sont lourdes pour ce genre d'infraction.
\end{quote}

\begin{quote}
\emph{The penalties are heavy for this kind of offence.}
\end{quote}

It can also mean approval, but mostly it is a penalty.

\begin{cousinweb}
\emph{sanction} is a learned borrowing from Latin \emph{sānctiō} \textquotedblleft{}a decree, a penalty clause\textquotedblright{} (EtymDB; Lewis \& Short), from \emph{sancīre} \textquotedblleft{}to make sacred, to decree\textquotedblright{}.
\end{cousinweb}

\begin{grammarlens}[title={What you've built}]
One more word for telling, explaining and describing at B1.
\end{grammarlens}

\subsection*{Your turn}
\begin{itemize}
\raggedright
  \item \emph{Say it:} \emph{la sanction}
  \item \emph{Say it:} \emph{la sanction}, once more
  \item \emph{Say it:} a sentence of your own with \emph{la sanction}
  \item \emph{Recall:} say the French for the offence, then the French for to transgress, then say \emph{la sanction} again
\end{itemize}

\begin{tcolorbox}[breakable,colback=teal!4,colframe=teal!35!black,title={Before you move on}]
What does la sanction mean? (\textquotedblleft{}The penalty, the sanction\textquotedblright{}.) Say it once more.
\end{tcolorbox}
\section[sanctionner]{sanctionner — to penalise, to sanction}

\label{lesson:FR-C382-sanctionner}



\begin{quote}
Before the new one: say the French for to transgress, then the French for the penalty.
\end{quote}

\subsection*{You'll want to know: sanctionner}
\textbf{sanctionner} — \textquotedblleft{}to penalise, to sanction\textquotedblright{}.

\begin{quote}
L'école a sanctionné les élèves qui fumaient dans les toilettes.
\end{quote}

\begin{quote}
\emph{The school penalised the pupils who were smoking in the toilets.}
\end{quote}

\emph{Être sanctionné} — to be penalised.

\begin{cousinweb}
\emph{sanctionner} is built on \emph{sanction}, from Latin \emph{sānctiō} \textquotedblleft{}a decree, a penalty clause\textquotedblright{}, from \emph{sancīre} \textquotedblleft{}to make sacred, to decree\textquotedblright{} (Lewis \& Short).
\end{cousinweb}

\begin{grammarlens}[title={What you've built}]
One more word for telling, explaining and describing at B1.
\end{grammarlens}

\subsection*{Your turn}
\begin{itemize}
\raggedright
  \item \emph{Say it:} \emph{sanctionner}
  \item \emph{Say it:} \emph{sanctionner}, once more
  \item \emph{Say it:} a sentence of your own with \emph{sanctionner}
  \item \emph{Recall:} say the French for to transgress, then the French for the penalty, then say \emph{sanctionner} again
\end{itemize}

\begin{tcolorbox}[breakable,colback=teal!4,colframe=teal!35!black,title={Before you move on}]
What does sanctionner mean? (\textquotedblleft{}To penalise, to sanction\textquotedblright{}.) Say it once more.
\end{tcolorbox}
